% Section 07 - Configuration académique
\section{Configuration académique des classes}

\subsection{Menu "Config des relevés"}

Ce menu permet de créer et gérer les configurations de classes nécessaires pour générer les relevés de notes.

\subsection{Structure hiérarchique}

Une configuration de classe est organisée en 3 niveaux :

\begin{verbatim}
Configuration de Classe
├── Semestre 1
│   ├── UE1 - Sciences Fondamentales (30 crédits)
│   │   ├── EC: Mathématiques (poids: 2, base: 20)
│   │   ├── EC: Physique (poids: 1.5, base: 20)
│   │   └── EC: Chimie (poids: 2, base: 20)
│   └── UE2 - Sciences Humaines (30 crédits)
│       └── EC: Anglais (poids: 1, base: 20)
├── Semestre 2
│   └── ...
└── Semestre 3
    └── ...
\end{verbatim}

\subsection{Créer une nouvelle configuration}

\subsubsection{Étape 1 : Informations générales}

\begin{enumerate}[leftmargin=*]
    \item Cliquez sur \bouton{Nouvelle Configuration}
    \item Remplissez :
    \begin{itemize}
        \item \champ{Nom de la configuration} : Ex. "Licence 1 Médecine - 2023-2024"
        \item \champ{Niveau} : L1, L2, L3, M1, M2, etc.
    \end{itemize}
    \item Cliquez sur \bouton{Créer}
\end{enumerate}

\subsubsection{Étape 2 : Ajouter des semestres}

\begin{enumerate}[leftmargin=*]
    \item Dans la configuration créée, cliquez sur \bouton{Ajouter un semestre}
    \item Remplissez :
    \begin{itemize}
        \item \champ{Nom du semestre} : "Semestre 1" ou "S1"
        \item \champ{Code} : S1, S2, S3, etc.
    \end{itemize}
    \item Répétez pour ajouter tous les semestres nécessaires
\end{enumerate}

\subsubsection{Étape 3 : Ajouter des UEs}

Pour chaque semestre :

\begin{enumerate}[leftmargin=*]
    \item Cliquez sur \bouton{Ajouter une UE}
    \item Remplissez :
    \begin{itemize}
        \item \champ{Nom de l'UE} : "UE1 - Sciences Fondamentales"
        \item \champ{Crédits ECTS} : Nombre de crédits (ex: 30)
    \end{itemize}
    \item Cliquez sur \bouton{Créer l'UE}
\end{enumerate}

\subsubsection{Étape 4 : Ajouter des ECs}

Pour chaque UE :

\begin{enumerate}[leftmargin=*]
    \item Cliquez sur \bouton{Ajouter un EC}
    \item Remplissez :
    \begin{itemize}
        \item \champ{Nom de l'EC} : "Mathématiques", "Physique", etc.
        \item \champ{Poids (coefficient)} : 1, 1.5, 2, etc.
        \item \champ{Base de notation} : 10, 20, ou 100
        \item \champ{A une session de rattrapage} : Cochez si applicable
    \end{itemize}
    \item Cliquez sur \bouton{Créer l'EC}
    \item Répétez pour tous les ECs de l'UE
\end{enumerate}

\begin{figure}[h]
    \centering
    % \includegraphics[width=\textwidth]{images/config_classe.png}
    \fbox{\parbox{0.9\textwidth}{\centering [CAPTURE D'ÉCRAN : INTERFACE DE CONFIGURATION]\\ \vspace{0.3cm}
    Arborescence : Config > Semestre > UE > EC\\
    Boutons "Ajouter" à chaque niveau\\
    Formulaires d'édition}}
    \caption{Interface de configuration académique}
\end{figure}

\subsection{Réorganiser les UEs}

Par défaut, les UEs s'affichent dans l'ordre de création. Pour les réorganiser :

\begin{enumerate}[leftmargin=*]
    \item Cliquez sur \bouton{Réorganiser les UEs}
    \item Utilisez le drag \& drop pour réordonner
    \item Ou utilisez les boutons ↑ et ↓
    \item Cliquez sur \bouton{Enregistrer l'ordre}
\end{enumerate}

\subsection{Personnaliser le thème des relevés}

Chaque configuration peut avoir son propre thème :

\begin{enumerate}[leftmargin=*]
    \item Dans la configuration, cliquez sur l'onglet \textbf{Thème}
    \item Configurez :
    \begin{itemize}
        \item Couleurs (primaire, secondaire, accent)
        \item Polices (famille, tailles)
        \item Mise en page
    \end{itemize}
    \item Prévisualisez en temps réel
    \item Enregistrez le thème
\end{enumerate}

\subsection{Exporter / Importer une configuration}

\subsubsection{Exporter}

\begin{enumerate}[leftmargin=*]
    \item Sélectionnez la configuration
    \item Cliquez sur \bouton{Exporter la configuration}
    \item Un fichier JSON est téléchargé : \texttt{config\_[Nom].json}
\end{enumerate}

\subsubsection{Importer}

\begin{enumerate}[leftmargin=*]
    \item Cliquez sur \bouton{Importer une configuration}
    \item Sélectionnez le fichier JSON
    \item La configuration est chargée et ajoutée à votre liste
\end{enumerate}

\astuce{Partagez vos configurations avec vos collègues via export/import pour standardiser les structures académiques.}

\subsection{Menu "Configurer les entêtes"}

Ce menu permet de définir les informations de l'établissement qui apparaîtront sur tous les documents.

\subsubsection{Informations de base}

\begin{itemize}[leftmargin=*]
    \item \champ{Nom de l'établissement}
    \item \champ{Adresse complète}
    \item \champ{Téléphone}
    \item \champ{Email}
    \item \champ{Site web}
\end{itemize}

\subsubsection{Logo}

\begin{enumerate}[leftmargin=*]
    \item Cliquez sur \bouton{Charger un logo}
    \item Sélectionnez une image PNG ou JPG
    \item Recommandations :
    \begin{itemize}
        \item Résolution : 800x200 pixels minimum
        \item Taille : Maximum 2 Mo
        \item Format : PNG avec transparence recommandé
    \end{itemize}
    \item Le logo s'affichera en haut de tous les relevés et attestations
\end{enumerate}

\subsection{Menu "Historique des docs"}

Ce menu permet de consulter l'historique de tous les documents générés.

\subsubsection{Informations affichées}

\begin{itemize}[leftmargin=*]
    \item Date et heure de génération
    \item Type de document (Relevé / Attestation)
    \item Nombre de documents générés
    \item Configuration utilisée
    \item Semestre(s) concerné(s)
    \item Format d'export utilisé
\end{itemize}

\subsubsection{Recherche et filtrage}

\begin{enumerate}[leftmargin=*]
    \item Utilisez la barre de recherche pour trouver une génération spécifique
    \item Filtrez par :
    \begin{itemize}
        \item Type de document
        \item Date (période)
        \item Configuration
    \end{itemize}
    \item Cliquez sur une entrée pour voir les détails complets
\end{enumerate}

\subsubsection{Export de l'historique}

\begin{enumerate}[leftmargin=*]
    \item Cliquez sur \bouton{Exporter l'historique}
    \item Choisissez le format : Excel ou PDF
    \item Un fichier récapitulatif est téléchargé
\end{enumerate}

\newpage
